% Consolidated guide-selftest05 from 97b62cc61307eba158a4498b55de3a5fdef3ff43; source fragments preserved.

% === BEGIN guide-selftest05-choices.tex blob a206fa1e6cad4bdad3af5daa0b6047849e50cb77 ===
% 第16批：扫描原书自测题5选择，共10题；题意独立转述。
\subsection{自测题5：气体动理论和热力学选择题}

\sourcetest{5.1.1}{115}{123}
\problemtext 封闭刚性容器中理想气体分子平均速率变为原来的2倍，问绝对温度和压强的倍率。
\workings 速率尺度 $\bar v\propto\sqrt T$，故 $T'=4T$；等容且粒子数不变，$p\propto T$，于是 $p'=4p$，选 \textbf{D}。
\crosscheck 题设速率是统计平均量；结论使用同种气体、理想气体、恒容前提。

\sourcetest{5.1.2}{115}{123}
\problemtext 活塞迅速外移而气体膨胀，判别气体是否经历平衡态以及能否以单条 $p$--$V$ 准静态曲线表述。
\workings 迅速过程通常存在非均匀压强和流动，系统内无单一瞬时热力学压强，非准静态平衡过程，不能以连续平衡态曲线准确表示，选 \textbf{C}。
\crosscheck 初态和充分弛豫后的末态可有明确 $p,V$，不等于中间全过程有平衡态曲线。

\sourcetest{5.1.3}{115}{123}
\problemtext 相同容积容器中的氧气与氦气初温初压相同；给氧气5J使其升高某温差，氦气升高同温差需多少热量？
\workings 同温同压同体积意味着粒子数相同。常温氧气刚性双原子，$C_V=5R/2$；氦气单原子，$C_V=3R/2$。等容下 $\boxed{Q_{\rm He}/Q_{\rm O_2}=3/5}$，故 $Q_{\rm He}=3$J，选 \textbf{C}。
\crosscheck 按理想气体、忽略振动态激发的常温近似；容器均刚性。

\sourcetest{5.1.4}{115}{123}
\problemtext 题5.1.2图中两个卡诺循环共用高低等温线，但第二个循环面积较大；比较净功与效率。
\workings 一个循环的净功是 $pV$ 图中围成的面积；第二循环面积变大故净功增加。卡诺效率仅与热源温度相关，$\eta=1-T_2/T_1$ 不变，选 \textbf{D}。
\crosscheck 题图改变循环体积尺度，未改变冷热源温度。

\sourcetest{5.1.5}{115}{123}
\problemtext 在温度为327°C和27°C的热源之间工作的理想可逆热机，最大效率是多少？
\workings $T_h=600.15$K、$T_c=300.15$K，$\eta_{\max}=1-T_c/T_h\approx\boxed{50\%}$，选 \textbf{B}。
\crosscheck 应换算成开尔文后才取比；按通常近似273K换算也得到50%。

\sourcetest{5.1.6}{115--116}{NaN}
\problemtext 理想气体从 $(V_0,T_0)$ 先绝热膨胀到 $2V_0$，等容升温回 $T_0$，再等温压缩回初态；判断循环净功。
\workings 令 $\gamma>1$。绝热末温 $T_a=T_0\,2^{1-\gamma}$，绝热功 $W_a=nRT_0[1-2^{1-\gamma}]/(\gamma-1)$；等容功0，等温压缩功 $W_c=-nRT_0\ln2$。因 $1-e^{-z}<z$（$z=(\gamma-1)\ln2$），$W_a+W_c<0$，选 \textbf{B}。
\crosscheck 一循环 $\Delta U=0$，净功为负意味着需净输入功。

\sourcetest{5.1.7}{116}{124}
\problemtext 一定量理想气体在真空中绝热自由膨胀，体积增大；问内能、熵变化。
\workings 无热量交换 $Q=0$，且向真空不作体积功 $W=0$，故 $\Delta U=0$、$T$ 不变；熵变 $\Delta S=nR\ln(V_2/V_1)>0$，选 \textbf{A}。
\crosscheck 过程不可逆，不能把实际膨胀误判为可逆绝热等熵过程。

\sourcetest{5.1.8}{116}{124}
\problemtext 理想气体只接触一个热源作等温膨胀，吸热全部转化为对外做功；这是否违反热力学定律？
\workings 对理想气体等温 $\Delta U=0$，由第一定律 $Q=W>0$；热力学第二定律的开尔文表述禁止的是\emph{循环}唯一效果，不禁止单次等温膨胀。两定律均不违反，选 \textbf{C}。
\crosscheck 若要回到初态必须另行进行复原步骤，不能将此过程声称为单热源循环热机。

\sourcetest{5.1.9}{116}{124}
\problemtext 刚性双原子氮气经可逆绝热压缩，压强增至2倍，求分子平均速率的倍率。
\workings $\gamma=7/5$，$T_2/T_1=(p_2/p_1)^{(\gamma-1)/\gamma}=2^{2/7}$；$\bar v_2/\bar v_1=\sqrt{T_2/T_1}=\boxed{2^{1/7}}$，选 \textbf{D}。
\crosscheck 若问分子均方根速率、平均速率或最概然速率，固定分子质量下倍率均为 $\sqrt T$。

\sourcetest{5.1.10}{116}{124}
\problemtext 理想气体绝热向真空自由膨胀，体积翻倍；比较初末温度与平均自由程。
\workings 自由膨胀理想气体 $\Delta U=0$，故 $T_1=T_2$；气体数密度减半，固定温度时 $\lambda\propto1/n$，故 $\lambda_2=2\lambda_1$。选 \textbf{B}。
\crosscheck 需要同种气体的有效碰撞直径不变；与可逆绝热降温不同。
% === END guide-selftest05-choices.tex ===

% === BEGIN guide-selftest05-blanks.tex blob a91ca00f32298644299b74444f3bef4bdff48614 ===
% 第16批：原书印刷第116—117页；源图题标注复核条件。
\subsection{自测题5：理想气体和热力学填空}

\sourcetest{5.2.1}{116}{124}
\problemtext 氢分子质量 $3.3\times10^{-24}$g，每秒有 $10^{23}$ 个分子以 $10^5$cm/s 速度、与壁面法线夹45°，在 $2.0$cm$^2$ 面积内完全弹性反弹；求压强。
\workings 单次法向冲量 $2mv\cos45^\circ$，$p=N(2mv\cos45^\circ)/(A\Delta t)$。代入 $m=3.3\times10^{-27}$kg、$v=10^3$m/s、$A=2.0\times10^{-4}$m$^2$、$N/\Delta t=10^{23}$s$^{-1}$，得 $\boxed{p\simeq2.33\times10^3\ {\rm Pa}}$。
\crosscheck 法线方向速度是 $v\cos45^\circ$；不能把单次碰壁动量变化写成 $mv$。

\sourcetest{5.2.2}{116}{124}
\problemtext 理想气体密度在等压升温、等温增压两种过程中分别怎样改变？
\workings 由 $\rho=pM/(RT)$，等压时 $\rho\propto1/T$，\emph{减小}；等温时 $\rho\propto p$，\emph{增大}。
\crosscheck 摄氏温度不可直接代入理想气体绝对温度比例。

\sourcetest{5.2.3}{116}{124}
\problemtext 氢气作为刚性双原子理想气体，升温1K时内能增加41.6J；求气体质量。
\workings $\Delta U=n(5R/2)\Delta T$，故 $n=41.6/(2.5\times8.31)\simeq2.00$mol。取氢气摩尔质量2g/mol，$\boxed{m\simeq4.00\ {\rm g}}$。
\crosscheck 采用常温转动自由度已激发、振动自由度未激发的刚性双原子近似。

\sourcetest{5.2.4}{117}{125}
\problemtext 定量理想气体先等容加热至原来绝对温度4倍，后等温膨胀至原来体积2倍；求分子平均撞击器壁频率倍率。
\workings 撞壁频率尺度 $z\propto n\bar v\propto (N/V)\sqrt T$，终态体积加倍、温度变4倍，故 $z_2/z_1=(1/2)\cdot2=\boxed{1}$。
\crosscheck 在固定截面积或单位面积的撞壁频率意义下成立；并非单个分子对随动器壁的碰撞计数。

\sourcetest{5.2.5}{117}{125}
\problemtext 在 $p$--$V$ 图中，分别用何种几何对象表征平衡态、准静态过程和循环过程？
\workings 一个平衡态用\emph{一个点}；一段准静态过程用\emph{一条曲线}；闭合准静态循环用\emph{一条闭合曲线}。
\crosscheck 非准静态过程的中间态通常无法赋予唯一的系统压强。

\sourcetest{5.2.6}{117}{125}
\problemtext 题5.2.2图中理想气体沿 $a\to b$ 等容升压、再 $b\to c$ 等压膨胀；判断气体吸热 $Q$ 与内能变化 $\Delta U$ 的符号。
\workings 两段均使 $T\propto pV$ 上升，$\Delta U>0$；第二段气体对外做功 $W>0$、首段 $W=0$，由 $Q=\Delta U+W$ 知 $\boxed{Q>0,\ \Delta U>0}$。
\crosscheck 初末温度不同，图中经过原点的虚线仅为状态标尺，不是等温线。

\sourcetest{5.2.7}{117}{125}
\problemtext 1kg水由30°C升至80°C，定压比热约4.2J/(g K)，同期电流做功 $4.2\times10^5$J；求水从外界通过热传递获得的热量（正负号）。
\workings 水内能增加近似 $\Delta U=1000\cdot4.2\cdot50=2.1\times10^5$J；电功输入 $W_e=4.2\times10^5$J，故传热 $Q=\Delta U-W_e=\boxed{-2.1\times10^5\ {\rm J}}$。
\crosscheck 负号表示水向环境净散热；要区分传热与电能输入。

\sourcetest{5.2.8}{117}{125}
\problemtext 1mol刚性双原子理想气体在等压升温1K中，内能增加20.78J；求对外做功和吸热。
\workings 由理想气体等压功 $W=nR\Delta T=\boxed{8.31\ {\rm J}}$；第一定律得 $Q=\Delta U+W=20.78+8.31=\boxed{29.09\ {\rm J}}$。
\crosscheck 该20.78J对应五自由度的 $5R/2$ 热容，和题给内能变化一致。

\sourcetest{5.2.9}{117}{125}
\problemtext 3mol单原子理想气体，初态 $p_1=6$atm、$T=500$K，等温膨胀至 $p_2=3$atm；求吸热。
\workings 等温理想气体 $\Delta U=0$、$V_2/V_1=p_1/p_2=2$，故 $\boxed{Q=W=nRT\ln2=3(8.31)(500)\ln2\simeq8.64\times10^3\ {\rm J}}$。
\crosscheck 结果与单原子自由度无关，因理想气体等温内能不变。
% === END guide-selftest05-blanks.tex ===

% === BEGIN guide-selftest05-solutions.tex blob bffecbe79d42dab7ad71ce90fa24e3df8efaa5a5 ===
% 第16批：独立推导、自查；依赖源图的题标明判读前提。
\subsection{自测题5：统计物理、热力学计算与证明}

\sourcetest{5.3.1}{117}{125}
\problemtext 体积1m$^3$ 中有 $N_1=1.0\times10^{25}$ 个氧分子和 $N_2=4.0\times10^{25}$ 个氮分子，混合气体压强 $2.76\times10^5$Pa；求每分子平均平动动能及温度。
\workings 总粒子数 $N=5.0\times10^{25}$，由 $pV=Nk_BT$ 得 $k_BT=pV/N=5.52\times10^{-21}$J。平均平动能 $\boxed{\bar\varepsilon_{\rm tr}=3k_BT/2=8.28\times10^{-21}\ {\rm J}}$；取 $k_B=1.380649\times10^{-23}$J/K，$\boxed{T\simeq399.8\ {\rm K}}$。
\crosscheck 混合气体平衡态两种分子的每粒子平均平动能相同，与分子质量无关。

\sourcetest{5.3.2}{117}{125}
\problemtext 氢核（质子）温度达到 $10^8$K，以非相对论理想气体近似求方均根速率与平均平动能（eV）。
\workings $\boxed{v_{\rm rms}=\sqrt{3k_BT/m_p}\simeq1.57\times10^6\ {\rm m/s}}$，取 $m_p=1.673\times10^{-27}$kg。$\bar\varepsilon=3k_BT/2=2.071\times10^{-15}$J；以 $1\,{\rm eV}=1.602176634\times10^{-19}$J，$\boxed{\bar\varepsilon\simeq1.29\times10^4\,{\rm eV}=12.9\,{\rm keV}}$。
\crosscheck $v_{\rm rms}\ll c$，非相对论温度速度估计自洽。

\sourcetest{5.3.3}{117}{125}
\problemtext 15°C、压强0.76mHg时氩气和氖气平均自由程分别为 $6.7\times10^{-8}$m、$13.2\times10^{-8}$m；求有效碰撞直径比 $d_{\rm Ne}/d_{\rm Ar}$；氩在20°C、0.15mHg的平均自由程。
\workings 硬球气体 $\lambda=k_BT/(\sqrt2\pi d^2p)$，同态下 $\boxed{d_{\rm Ne}/d_{\rm Ar}=\sqrt{\lambda_{\rm Ar}/\lambda_{\rm Ne}}=\sqrt{6.7/13.2}\simeq0.712}$。同种氩气改变状态时 $\lambda'_{{\rm Ar}}=\lambda_{\rm Ar}(T_2/T_1)(p_1/p_2)$；$T_1=288.15$K，$T_2=293.15$K，故 $\boxed{\lambda'_{{\rm Ar}}\simeq6.7\times10^{-8}(293.15/288.15)(0.76/0.15)\simeq3.45\times10^{-7}\ {\rm m}}$。
\crosscheck 两压强单位均为米汞柱，可直接取比；有效直径是简化硬球碰撞参数。

\sourcetest{5.3.4}{117--118}{NaN}
\problemtext 2mol氢气初态温度27°C、体积20L，先等压膨胀使体积加倍，再可逆绝热膨胀至温度回到初温；求 $pV$ 过程、吸热、内能变化与总功。
\workings 以刚性双原子 $\gamma=7/5$ 计算。$T_1\simeq300$K、$V_1=20$L；等压末态 $V_2=40$L、$T_2=600$K；随后绝热末体积 $V_3=V_2(T_2/T_3)^{1/(\gamma-1)}=40\cdot2^{5/2}\simeq226.3$L。等压吸热 $Q_{12}=n(7R/2)(T_2-T_1)=2(7R/2)(300)\simeq\boxed{1.75\times10^4\ {\rm J}}$，绝热段 $Q_{23}=0$，总 $\Delta U=nC_V(T_3-T_1)=\boxed{0}$；故全过程 $W=Q-\Delta U=\boxed{1.75\times10^4\ {\rm J}}$。
\crosscheck 这是开放路径而非热机循环，$T_3=T_1$ 不代表 $V_3=V_1$。

\sourcetest{5.3.5}{118}{126}
\problemtext 氮气初始1atm、300K，经可逆绝热压缩至32atm；求末温及末态粒子数密度。
\workings 取 $\gamma=7/5$；$T_2=T_1(p_2/p_1)^{(\gamma-1)/\gamma}=300\cdot32^{2/7}\simeq\boxed{807.6\ {\rm K}}$。数密度 $n_N=p_2/(k_BT_2)$，用 $1$atm$=101325$Pa，$\boxed{n_N\simeq(32\cdot101325)/(1.380649\times10^{-23}\cdot807.6)\simeq2.91\times10^{26}\ {\rm m^{-3}}}$。
\crosscheck 这里 $n_N$ 是每立方米分子数，并非摩尔数；绝热指数假设分子振动自由度未激发。

\sourcetest{5.4.1}{118}{126}
\problemtext 证明两恒温热源间可逆卡诺热机效率 $\eta=1-T_c/T_h$。
\workings 卡诺循环由温度 $T_h,T_c$ 的两段可逆等温与两段可逆绝热组成。可逆绝热段 $\Delta S=0$，故两等温段气体熵变大小相等：$Q_h/T_h=-Q_c/T_c$，其中 $Q_h>0$、$Q_c<0$。整循环 $\Delta U=0$，输出功 $W=Q_h+Q_c$，所以 $\boxed{\eta=W/Q_h=1+Q_c/Q_h=1-T_c/T_h}$。
\crosscheck 只对两个恒温热源间可逆卡诺机成立；不可逆实际热机效率通常更低。
% === END guide-selftest05-solutions.tex ===
